% !TeX root = ../深度学习与大语言模型技术原理.tex
% 由主文档通过 \input 载入；请编译主文档。
\setcounter{chapter}{24} % 保留原书章号。
\chapter{推测解码与预填充解码分离}
\label{ch:speculation}

\section{草稿与目标模型的并行验证}

推测解码（speculative decoding）让较便宜的草稿模型先提出若干候选词元，目标模型再一次处理这些候选前缀，计算各位置的目标分布。因果掩码允许目标模型并行求得所有候选条件概率，虽然最终只能接受到第一个拒绝位置。其目标是减少昂贵目标模型的串行调用次数，而非减少所有模型的总运算量\cite{speculative}。

设同一历史下草稿分布为 $q$，目标分布为 $p$，先采样 $Y\sim q$，以 $a(Y)=\min(1,p(Y)/q(Y))$ 接受。被拒绝时改从剩余分布
\begin{equation}
r(v)=\frac{[p(v)-q(v)]_+}{\sum_u[p(u)-q(u)]_+}
\label{eq:speculative-residual}
\end{equation}
采样。$q(Y)>0$ 对实际草稿样本自动成立；当分母为零时 $p=q$，不会发生拒绝。

\begin{proposition}
\label{prop:speculative-exact}
在相同有限词表上精确使用上述接受与残差采样规则，最终词元分布为 $p$。
\end{proposition}
\begin{proof}
词元 $v$ 通过草稿接受得到的概率为 $q(v)\min(1,p(v)/q(v))=\min(p(v),q(v))$，$q(v)=0$ 时该项为零。总拒绝概率为
\[
1-\sum_u\min(p(u),q(u))=\sum_u[p(u)-q(u)]_+.
\]
乘以式~\eqref{eq:speculative-residual}，拒绝后输出 $v$ 的概率为 $[p(v)-q(v)]_+$。两项相加恰为 $p(v)$。
\end{proof}

对多词元草稿逐位置应用该规则，在已接受历史上继续检查，遇到拒绝则采样纠正词元并丢弃后续草稿；全部接受时可从最后的目标分布再采一个词元。要保持目标分布，温度、top-$p$ 与其他约束应先一致地纳入 $p,q$，不能混用处理前后的概率。贪心匹配验证则是另一种确定性约定。

\section{接受率与加速条件}

若每次草稿长度为 $K$，每步草稿成本为 $t_d$，目标验证成本为 $t_v(K)$，额外成本为 $t_o$，每轮平均提交 $\mathbb E[A]$ 个词元，则平均每词元成本约为 $(Kt_d+t_v(K)+t_o)/\mathbb E[A]$。只有低于普通目标解码成本时才有净加速。

在每个条件接受率都近似常数 $a$、忽略提前结束且包含全接受后的额外词元时，$\mathbb E[A]=1+a+\cdots+a^K$。实际接受率随任务和前缀变化，验证成本也会随 $K$ 与批量增长，所以更长草稿不必更快。高并发下目标模型本已充分利用算力，额外草稿可能削弱收益。

缓存回滚也是正确性的一部分：拒绝位置之后的草稿 KV 不能作为已确认历史保留。分支草稿树需要特定掩码和索引，不能把不同分支当成一条连续句子执行。

\section{EAGLE、多词元预测与词串草稿}

EAGLE 使用目标模型特征构建轻量草稿，利用词元信息处理特征预测的不确定性；不同版本会改变草稿结构与验证调度\cite{eagle}。其意义是把草稿器与目标模型内部表示结合，而非所有推测方法都必须加载一个独立完整小模型。

多词元预测（multi-token prediction, MTP）在训练时增加未来多个位置的预测头，例如共享表示 $h_t$ 分别预测 $x_{t+1},\ldots,x_{t+K}$，目标可写为各头交叉熵的加权和\cite{mtp}。这些边际预测不自动构成正确的联合生成分布；用于推测时仍需合适草稿组织与目标验证。它也可以作为训练辅助目标，不能与推测解码完全等同。

词串草稿（n-gram speculation）从提示或已生成文本的重复片段提出后续词元，额外模型成本低，适合重复度较高的任务。草稿来源可以不同，分布保持与否取决于完整验证算法，不取决于方法是否被称为“推测”。

\section{预填充解码分离与服务框架}

预填充解码分离（prefill--decode disaggregation）把两阶段放到不同资源池，预填充完成后传输 KV 给解码实例。这样可分别选择并行方式和扩容比例，降低两类计算互相干扰；DistServe 给出了围绕服务质量约束的系统研究\cite{distserve}。代价是缓存传输、跨池队列和状态协调。

若需传输 $M_{\mathrm{KV}}$ 字节、有效链路带宽为 $\mathcal B$，传输时间至少受 $M_{\mathrm{KV}}/\mathcal B$ 量级约束，此外还有启动、格式转换和排队。长提示产生大缓存时，网络可能抵消分离收益。应比较满足 TTFT 与 ITL 约束的有效吞吐量（goodput），而不只比较理论总吞吐。

vLLM、SGLang 和 TensorRT-LLM 分别提供不同模型执行、缓存与内核组织，具体特性和兼容矩阵随版本变化\cite{vllm,sglang,trt}。部署时固定模型、模板、内核与版本，测量冷启动、稳定负载、突发流量和过载恢复。服务层还要处理队列上限、超时、取消、流式输出和健康检查；算法优化只有在完整服务链路中保留收益，才算完成部署验证。

